% ===== Document Preamble =====
\documentclass[11pt]{article}
\usepackage[margin=0.5in]{geometry}
\usepackage{amsmath,amssymb}
\usepackage{xcolor}
\usepackage{enumitem}
\usepackage{needspace}
\usepackage{comment}

% ===== Colors =====
\definecolor{AnsRed}{RGB}{170,0,0}
\definecolor{WrongBlue}{RGB}{0,0,160}

% ===== Choice list styling (tight, consistent) =====
\setlist[enumerate]{%
  leftmargin=1.5em,
  itemsep=0.25em,
  topsep=0.25em,
  label=\textbf{\Alph*.}
}

% ===== Style Wrappers =====
\newcommand{\QuestionStyle}[1]{{\large\bfseries\boldmath #1}}
\newcommand{\AnswerStyle}[1]{{\large\bfseries\boldmath #1}}
\newcommand{\ReasonStyle}[1]{{\normalsize #1}}

% ===== Show / Hide Reasoning Toggle =====
\newif\ifshowreason
\showreasontrue      % Instructor version
% \showreasonfalse   % Student version

% ===== Choice Item Macros (Answer + Reasoning) =====
\newcommand{\Correct}[2]{%
  \item \AnswerStyle{\textcolor{AnsRed}{#1}}%
  \ifshowreason\\[-0.2em]
    \ReasonStyle{\textcolor{AnsRed}{#2}}%
  \fi
}

\newcommand{\Wrong}[2]{%
  \item \AnswerStyle{\textcolor{WrongBlue}{#1}}%
  \ifshowreason\\[-0.2em]
    \ReasonStyle{\textcolor{WrongBlue}{#2}}%
  \fi
}

% ===== Question Block Macro =====
\newcommand{\QBlock}[3]{%
  \Needspace{8\baselineskip}%
  \noindent
  \begin{minipage}{\linewidth}
    \QuestionStyle{#1.~#2}

    \vspace{0.25em}
    \begin{enumerate}
      #3
    \end{enumerate}

    \vspace{0.25em}
    \hrule height 0.4pt
  \end{minipage}
  \vspace{0.75em}
}

% ===== Document Start =====

\begin{document}

\begin{center}
{\huge \textbf{Physics Mechanics}}\\[0.25em]
{\Large Motion and Trajectory}}\\[0.25em]
{\large Multiple choice questions generated below. Answers in RED and incorrect answers in BLUE.}
\end{center}

\vspace{0.75em}

\QBlock{Q1}{A grasshopper jumps off the ground with a speed $u_0$ 
 at an angle \(\theta\). What is the speed of the grasshopper when it reaches the highest point of its trajectory?}{
  \Wrong{$u_0$.}{At maximum height, only the horizontal component 
  of velocity remains, not the initial speed.}
  \Correct{$u_0 \cos \theta$.}{At the maximum height, the upward 
  velocity is zero, leaving only the horizontal component, which 
  is $u_0 \cos \theta$.}
  \Wrong{$u_0 \sin \theta$.}{The vertical component of the velocity 
  becomes zero at the maximum height.}
  \Wrong{0.}{The total speed is not zero due to the presence of the 
  horizontal component.}
}

\QBlock{Q2}{A drone is moving horizontally with an initial velocity of 
\(\vec{v_0} = (25.5 \text{ m/s})\hat{i}\). It experiences an 
acceleration given by \(\vec{a} = (1.75 \text{ m/s}^2)\hat{i} + (3.28 \text{ m/s}^2)\hat{j}\) when its boosters 
activate. Which of the following statements is true about the drone's motion?}{
  \Wrong{The x-component of the drone's velocity remains constant.}{The x-component changes due to the 
x-component of the acceleration.}
  \Wrong{The y-component of the drone's velocity remains constant.}{The y-component changes due to the 
y-component of the acceleration.}
  \Correct{Both the direction of motion and speed change.}{The acceleration in both the x and y directions affects both 
the speed and the direction of motion.}
  \Wrong{None of the above is correct.}{The motion and speed change due to the given acceleration, 
confirming option C.}
}

\QBlock{Q3}{A projectile is launched from a tower with an initial speed of 10.2 m/s 
at an angle of 45$^\circ$ above the horizontal. Calculate the time it takes for the 
projectile to reach its maximum height.}{
  \Wrong{0.50 seconds.}{This option results from incorrectly calculating 
the time by using the wrong vertical component.}
  \Correct{0.72 seconds.}{The time to reach maximum height is calculated using 
$t = \frac{u_0 \sin(\theta)}{g}$. Here, $u_0 = 10.2 \text{ m/s}$ and $\theta = 45^\circ$, so 
$t = \frac{10.2 \times \sin(45^\circ)}{9.81} \approx 0.72 \text{ seconds}$.}
  \Wrong{0.90 seconds.}{This value is based on doubling the time for a 
different component of motion.}
  \Wrong{1 second.}{Assumes incorrect independence of time and angle.}
}

\end{document}